\subsection{Discussion}
\label{sec:discussion}

This section interprets the formal results from Section~\ref{sec:theoretical_properties}, showing how the measure relationships and scope-structure bounds translate into actionable diagnostic information. All three measures share identical decision-problem complexity (Theorems~\ref{thm:ur_complexity}--\ref{thm:gs_complexity}), so the choice of which profiles to compute should be driven by diagnostic needs rather than computational cost. The practical time-complexity differences documented in Table~\ref{tab:complexity} remain relevant for explicit enumeration approaches, where the Goal Sequencing Conflict Profile requires additional local reachability computations. The section then discusses normalization considerations for practical deployment.

\subsubsection{Diagnostic Interpretation}

The measure relationships established in Proposition~\ref{prop:measure_relationships} reflect increasing diagnostic specificity. Unreachability identifies goals that cannot be achieved at all; mutex identifies achievable goals that cannot coexist; sequencing identifies mutex goals where the exclusion is irreversible. Figure~\ref{fig:measure_relationships} illustrates this containment structure. Two cases are particularly informative:

\begin{itemize}
    \item \textbf{Mutex without sequencing} (e.g., Light Switch): Goals exclude
          each other but actions are reversible, so the conflict is symmetric
          and neither goal is more responsible for unsolvability.
    \item \textbf{Mutex with sequencing} (e.g., Trust and Travel): Goals exclude
          each other \emph{and} achieving specific goals creates permanent dead
          states. The sequencing measure identifies which goal achievement is
          responsible, providing actionable information for adding reversal
          operators.
\end{itemize}

State space analysis (unreachability, mutex) determines what states exist, while trajectory analysis (sequencing) determines what orderings of goal achievements are possible. For irreversible problems, all three measures report non-zero values, with sequencing providing directional information about which goal achievements lead to dead states.

The scope-structure bounds (Proposition~\ref{prop:scope_structure_bounds}) enable further diagnostic interpretation through ratio analysis:

\begin{itemize}
    \item \textbf{Unreachability depth}: When $\mathcal{I}^{\text{scope}}_{\text{UR}} > 0$, the ratio $(\mathcal{I}^{\text{struct}}_{\text{UR}} - \mathcal{I}^{\text{scope}}_{\text{UR}}) / \mathcal{I}^{\text{scope}}_{\text{UR}}$ indicates cascading dependency depth. A ratio of 0 means goals are directly unreachable, with no intermediate propositions blocked. Larger ratios indicate deeper precondition chains.

    \item \textbf{Mutex density}: If $\mathcal{I}^{\text{struct}}_{\text{MX}} = \binom{\mathcal{I}^{\text{scope}}_{\text{MX}}}{2}$, the mutex goals form a complete clique where every pair is mutually exclusive, as in Traffic Light. Ratios below 1 indicate sparser conflict structures with isolated mutex pairs.

    \item \textbf{Sequencing symmetry}: If $\mathcal{I}^{\text{struct}}_{\text{GS}} = \mathcal{I}^{\text{scope}}_{\text{GS}} \times (\mathcal{I}^{\text{scope}}_{\text{GS}} - 1)$, all sequencing conflicts are bidirectional: achieving any involved goal blocks all others, and vice versa. Ratios near 0.5 suggest predominantly unidirectional conflicts. For goal sets with three or more conflict participants, lower ratios indicate single-source dominant patterns where one goal blocks many others but not conversely.
\end{itemize}

Table~\ref{tab:measure_values} demonstrates these relationships across the test scenarios. For problems exhibiting only unreachability, only $\mathcal{I}_{\text{UR}}$ yields non-zero values. For reversible mutex problems, only $\mathcal{I}_{\text{MX}}$ yields non-zero values. Irreversible problems yield non-zero values for both $\mathcal{I}_{\text{MX}}$ and $\mathcal{I}_{\text{GS}}$, consistent with Proposition~\ref{prop:measure_relationships}.

\begin{table}[ht]
    \centering
    \begin{tabular}{l cc cc cc}
        \toprule
        Scenario         & \multicolumn{2}{c}{$\mathcal{I}_{\text{UR}}$} & \multicolumn{2}{c}{$\mathcal{I}_{\text{MX}}$} & \multicolumn{2}{c}{$\mathcal{I}_{\text{GS}}$}                           \\
        \cmidrule(lr){2-3}\cmidrule(lr){4-5}\cmidrule(lr){6-7}
                         & scope                                         & struct                                        & scope                                         & struct & scope & struct \\
        \midrule
        Locked Door      & 1                                             & 3                                             & 0                                             & 0      & 0     & 0      \\
        Bank Vault       & 1                                             & 7                                             & 0                                             & 0      & 0     & 0      \\
        Light Switch     & 0                                             & 0                                             & 2                                             & 1      & 0     & 0      \\
        Traffic Light    & 0                                             & 0                                             & 3                                             & 3      & 0     & 0      \\
        Trust and Travel & 0                                             & 0                                             & 2                                             & 1      & 2     & 1      \\
        Rival Alliances  & 0                                             & 0                                             & 2                                             & 1      & 2     & 2      \\
        \bottomrule
    \end{tabular}
    \caption{Measure values across test scenarios.}
    \label{tab:measure_values}
\end{table}

The complete diagnostic profile for a planning problem is the 6-tuple:
\begin{equation*}
    (\mathcal{I}^{\text{scope}}_{\text{UR}}(\Pi, h), \mathcal{I}^{\text{struct}}_{\text{UR}}(\Pi, h), \mathcal{I}^{\text{scope}}_{\text{MX}}(\Pi, h), \mathcal{I}^{\text{struct}}_{\text{MX}}(\Pi, h), \mathcal{I}^{\text{scope}}_{\text{GS}}(\Pi, h), \mathcal{I}^{\text{struct}}_{\text{GS}}(\Pi, h))
\end{equation*}
Table~\ref{tab:measure_values} lists the converged profiles $\text{Profile}(\Pi)$ for all test scenarios.

Each scenario in Table~\ref{tab:measure_values} exhibits a characteristic profile signature. Locked Door illustrates dependency-structure failure with depth 2, where one unreachable goal proposition appears among three unreachable propositions in total. Bank Vault shares the same scope value at depth 6, reflecting three converging precondition chains. The two reversible-mutex scenarios differ in clique completeness: Light Switch yields a single mutex pair, while Traffic Light closes a complete 3-clique with all $\binom{3}{2} = 3$ pairs mutex. Trust and Travel and Rival Alliances share mutex structure ($\mathcal{I}^{\text{struct}}_{\text{MX}} = 1$) but separate by sequencing direction: Trust and Travel exhibits unidirectional sequencing ($\mathcal{I}^{\text{struct}}_{\text{GS}} = 1$) and Rival Alliances exhibits bidirectional sequencing ($\mathcal{I}^{\text{struct}}_{\text{GS}} = 2$). This contrast illustrates how trajectory analysis (sequencing) provides finer-grained diagnostic information beyond state space analysis (unreachability, mutex). These interpretive observations rely on the absolute counts produced by the measures, raising the question of whether such counts should be normalized for cross-problem comparison.

\subsubsection{Normalization Considerations}

The three proposals presented above report absolute inconsistency values: counts of unreachable propositions, mutex pairs, and sequencing conflict pairs. A practical question is whether these measures should be normalized relative to problem characteristics such as the total number of goal propositions or the size of the proposition space.

Besnard and Grant~\cite{Besnard2020} develop relative inconsistency measures for propositional knowledge bases, demonstrating both advantages and limitations of normalization approaches. Relative measures can facilitate comparison across problems of different sizes, potentially providing a notion of percentage inconsistency analogous to test coverage metrics in software engineering.

However, several considerations suggest that absolute measures may be more appropriate for planning diagnostics. First, the severity of unsolvability is not necessarily proportional to problem size. A planning problem with 100 goal propositions where only 2 conflict may be easier to repair than a problem with 3 goal propositions where all 3 form a complete conflict graph. Second, normalization requires choosing an appropriate denominator from candidates such as the total goal-proposition count, the proposition-space size, or the maximum possible conflict count, and each choice encodes a different assumption about problem structure. Third, the proposed measures already provide multiple values that capture relative information: the ratio between unreachable goal propositions and total unreachable propositions indicates dependency depth, while the ratio between mutex goals and mutex pairs (or conflicting goals and sequencing pairs) reveals conflict structure.

For the scope of this thesis, the focus remains on absolute measures that report objective structural properties. Future work could calibrate normalization schemes against the \ac{IPC}~2016 benchmark distribution to compare inconsistency trends across problem generator parameters or to establish thresholds for automated problem classification.

Having defined the measures and analyzed their postulates, relationships, and complexity, Section~\ref{sec:implementation} turns to their computation, presenting the \ac{ASP}-based pipeline that realizes the parameterized profiles on standard \ac{PDDL} input.
